\section{Módulo 2: Planificación de Corto Plazo y MRP                                                        9}

\subsection{De lo agregado a lo desagregado . . . . . . . . . . . . . . . . . . . . . . . . . . . .       9}

\subsection{La mecánica de las tablas MRP . . . . . . . . . . . . . . . . . . . . . . . . . . . .         9}

\subsection{Lotificación: cómo decidir el tamaño del lote . . . . . . . . . . . . . . . . . . . . .      10}

\subsection{Comparación de los sistemas (ejemplo del curso) . . . . . . . . . . . . . . . . . .          11}

\subsection{MRP de ciclo cerrado y métricas TPM/OEE . . . . . . . . . . . . . . . . . . . .              11}

\subsection{V/F frecuentes de Planificación y MRP . . . . . . . . . . . . . . . . . . . . . . .          11}

\subsection{MRP: BOM y Silver-Meal (I2 2025) . . . . . . . . . . . . . . . . . . . . . . . . .           12}
\section{Módulo 2: Planificación de Corto Plazo y MRP}

\subsection{De lo agregado a lo desagregado}

\begin{definicion}{La cadena de desagregación}


 Planif. Agregada → MPS → MRP → Programación diaria

   MPS (Plan Maestro de Producción): traduce el plan agregado en productos es-
   pecíficos, cantidades específicas, períodos específicos. Responde “¿cuántas unidades del
   producto X produzco en la semana Y?”. No dice qué materiales se necesitan.
   MRP (Material Requirements Planning): explota la demanda hacia atrás en
   el tiempo, considerando los plazos de entrega, para saber cuándo lanzar órdenes.
   Desarrollado por J. Orlicky en IBM.

   En el corto plazo: períodos más cortos, productos a nivel individual, demanda menos
   incierta, y aparecen complejidades como los setups.

\end{definicion}

\begin{definicion}{Insumos que necesita el MRP}

1. BOM (Bill of Materials): la “receta” y estructura jerárquica del producto.
2. Plazo de entrega (L, lead time): tiempo entre emitir la orden y recibirla.
\end{definicion}
\section{Inventario disponible (OH, on hand ) de cada componente al comienzo.}

\subsection{La mecánica de las tablas MRP}

\begin{formulas}{Filas de la tabla MRP y su lógica}

   Para cada ítem se llena una tabla por semana con:

   GR (requerimiento bruto): demanda proveniente del nivel superior.
   OH (inventario disponible): stock al cierre del período. OHt = OHt−1 + recepcionest −
   GRt .
   NR (requerimiento neto) = GRt − OHt−1 − en tránsitot (si es > 0).
   POR (liberación de orden planificada): orden desfasada L períodos hacia atrás.

   Regla de oro: la orden del padre se convierte en GR del hijo, multiplicada por la cantidad
   de la receta, en la semana en que el padre la libera (POR), no en que la recibe.

\end{formulas}

\begin{trampa}{El desfase del Lead Time es donde todos se equivocan}

   Si una parte tiene L = 3 y necesito recibir 80 unidades en la semana 4 (NR), debo liberar
   la orden (POR) en la semana 1. El GR del componente hijo aparece en la semana
   1 (cuando el padre libera), no en la semana 4. Confundir “recepción” con “liberación”
   arrastra el error a todo el árbol BOM.






\end{trampa}

\subsection{Lotificación: cómo decidir el tamaño del lote}

\begin{definicion}{Las técnicas de lotificación}

   Todas equilibran costo de ordenar/setup vs. costo de mantener inventario:

   Lote a Lote (L4L): produces exactamente lo necesario cada período. Cero inventario,
   ideal para ítems caros/perecederos; pero muchos setups.
   Cantidad Fija de Pedido: siempre la misma cantidad. Simple, pero genera inventario
   innecesario o insuficiente.
   Período Fijo: ordenar cada N períodos. Simple, pero inventario variable.
   EOQ (cantidad económica de pedido): lote fijo óptimo bajo demanda constante; rara
   en MRP.
   Costo Total Mínimo: ajusta el lote buscando igualar el costo de ordenar con el de
   inventario acumulado.
   Costo Unitario Mínimo: divide el costo total entre las unidades del lote y minimiza
   ese cociente.
   Silver-Meal: heurística que minimiza el costo promedio por período.
   Wagner-Whitin (WW): el óptimo matemático exacto (prog. dinámica), pero
   costoso y poco intuitivo.

\end{definicion}

\begin{formulas}{Algoritmo Silver-Meal (paso a paso)}

   Se busca el lote que minimiza el costo promedio por período cubierto. Para un lote que
   arranca en el período 1 y cubre hasta el período k:
Pk
   S+  t=1 h · (t − 1) · Dt
   CT (k) =
k
   donde S = costo de setup, h = costo de mantener por unidad-período, Dt = demanda
   del período t. Regla de parada: se sigue agregando períodos al lote mientras CT (k)
   decrezca; en cuanto CT (k) > CT (k − 1), se cierra el lote en k − 1 y se reinicia el conteo
   desde k.

\end{formulas}

\begin{trampa}{Silver-Meal: no empieces en períodos con demanda cero}

   Si los primeros períodos tienen demanda = 0, sáltate directamente al primer período
   con demanda positiva y arranca Silver-Meal desde ahí. Q = 0 en los períodos con
   demanda cero es trivial: no hay nada que ordenar. Empezar ahí genera comparaciones
   degeneradas (C(1) = 0) que confunden sin aportar información.

\end{trampa}

\begin{tipprueba}{¿Cuándo usar cada método en la prueba?}

   Identificación: si el enunciado dice “use L4L” ⇒ produces exacto, OH siempre 0. Si dice
   “Silver-Meal” ⇒ tabla de costos promedio con regla de parada. Si dice “Wagner-Whitin”
   ⇒ matriz de horizonte (último período de producción vs. horizonte) y se elige el mínimo
   Zt . Qué aplicar: en Silver-Meal arma una mini-tabla de CT (k) por cada lote candidato;
   el inventario se valoriza por las semanas que “espera” (t − 1 períodos ×h).






\end{tipprueba}

\subsection{Comparación de los sistemas (ejemplo del curso)}

\begin{definicion}{Resultados comparativos (ejemplo de clase, 8 semanas)}

   Con S = $47, costo artículo $10, h = 0,5 % del valor:

   Método  Costo Total
   Lote a Lote (L4L)$376,00
   EOQ de Wilson$171,05
   Costo Unitario Mínimo$153,50
   Costo Total Mínimo   $140,05

   Conclusión clave: producir por lotes (no ajustarse a cada pedido) reduce costos. Por eso
   las empresas producen por lotes y no L4L puro.


\end{definicion}

\subsection{MRP de ciclo cerrado y métricas TPM/OEE}

\begin{definicion}{MRP de ciclo cerrado y OEE}

   El MRP de ciclo cerrado retroalimenta la planeación de capacidad: si el plan no es
   factible, se replantea el MPS y se repite el proceso período a período. En mantenimiento,
   el TPM (Mantenimiento Productivo Total) mide el rendimiento de los equipos mediante
   la OEE (Eficiencia Global del Equipo):

 D
OEE = Disponibilidad × Rendimiento × Calidad =
 A
   donde A=tiempo planificado, B=tiempo operativo, C=tiempo de funcionamiento,
   D=tiempo productivo, y Disp= B/A, Rend= C/B, Calidad= D/C.


\end{definicion}

\subsection{V/F frecuentes de Planificación y MRP}
  #RAfirmación y corrección
  1F“El MPS indica qué materiales y en qué cantidad comprar.” — Falso:
eso lo hace el MRP; el MPS solo dice cuántos productos terminados
producir.
  2V“El MRP explota la demanda hacia atrás considerando los lead times.”
  3F“L4L minimiza el número de setups.” — Falso: L4L minimiza el
inventario pero suele maximizar los setups.
  4V“Wagner-Whitin entrega la solución óptima exacta de lotificación.”
  5F“Agregar la demanda aumenta el error de pronóstico.” — Falso: la
agregación reduce el error relativo (los errores se compensan).
  6V“La estrategia Chase minimiza el inventario pero eleva los costos de
contratación/despido.”

\subsection{MRP — BOM, L4L y Silver-Meal (I2 2025, Pregunta 3)}

\begin{ejercicio}{Ensamble con BOM, L4L y Silver-Meal (I2 2025 / Ay. 8)}

\end{ejercicio}

\begin{tipprueba}{¿Cuándo usar este enfoque?}

   Identificación: producto final con BOM (1×A, 2×B, 1×C; A→A1+2A2;
   B→A1+B2), lead times, inventario inicial y orden en tránsito. Qué aplicar: (1)
   dibujar BOM, (2) tabla MRP del producto final con L4L (desfasar por L), (3)
   Silver-Meal para el ítem A1.

   i) BOM:

Producto Final

  A (1)B (2)  C (1)

 A1 (1)A2 (1)
   A1 (2)   B2 (1)

   ii) MRP del producto final (L4L). Demanda: sem 4=10, 5=50, 6=40, 7=60, 8=50.
   OH0 = 50; orden en tránsito 20 que llega en sem 2; L = 1 (ensamble).
Período   01   2 3   45 6  7 8
Req. Bruto 0   0 0  10   50 40 6050
Prod. Tránsito 0  20 0   00 0  0 0
Inv. Final 50 7070  60   10 0  0 0
Req. Neto  0   0 0   00 30 6050
Lote   0   0 0   00 30 6050
Orden Programada   0   0 0   0   30 60 500


Parte PF
Periodo   0  1  2  34 5   6 7 8
GR: Requerimiento Bruto 1050  406050
OH: Inventario Final  5070 70  70   6010   0 00
POR 3060  50

   Con L4L se produce exactamente el requerimiento neto, desfasado 1 semana.
   Explicación paso a paso de la mecánica de la tabla (MRP): Para entender cómo
   se obtiene cada número, aplicamos las siguientes fórmulas para cada semana t:

   Inventario al cierre (OHt ): OHt = OHt−1 + Recepciones Programadast +
   Recepciones Planificadast − GRt .
   Requerimiento  Neto (RNt ):   RNt=  máx{0, GRt − OHt−1   −
   Recepciones Programadast }.
   Liberación de Orden (P ORt ): Como L = 1, se desfasa una semana antes: P ORt−1 =
   RNt .








 Semana   Cálculo de Inventario (OH) y Requerimiento Neto (RN )

 Semana 1 OH1 = OH0 + Rec1 − GR1 = 50 + 0 − 0 = 50. Como OH1 ≥ 0 ⇒ RN1 = 0.
 Semana 2 OH2 = OH1 + Rec2 − GR2 = 50 + 20 − 0 = 70. Como OH2 ≥ 0 ⇒ RN2 = 0.
 Semana 3 OH3 = OH2 + Rec3 − GR3 = 70 + 0 − 0 = 70. Como OH3 ≥ 0 ⇒ RN3 = 0.
 Semana 4 OH4 = OH3 + Rec4 − GR4 = 70 + 0 − 10 = 60. Como OH4 ≥ 0 ⇒ RN4 = 0.
 Semana 5 OH5 = OH4 + Rec5 − GR5 = 60 + 0 − 50 = 10. Como OH5 ≥ 0 ⇒ RN5 = 0.
 Semana 6 El stock inicial es OH5 = 10 y la demanda es GR6 = 40.
  Nos faltan: RN6 = 40 − 10 = 30 unidades. Programamos recibir 30 ⇒ OH6 = 10 + 30 − 40 = 0.
  Con lead time L = 1, se libera una semana antes: POR5 = 30.
 Semana 7 El stock inicial es OH6 = 0 y la demanda es GR7 = 60.
  Nos faltan: RN7 = 60 − 0 = 60 unidades. Programamos recibir 60 ⇒ OH7 = 0 + 60 − 60 = 0.
  Con lead time L = 1, se libera una semana antes: POR6 = 60.
 Semana 8 El stock inicial es OH7 = 0 y la demanda es GR8 = 50.
  Nos faltan: RN8 = 50 − 0 = 50 unidades. Programamos recibir 50 ⇒ OH8 = 0 + 50 − 50 = 0.
  Con lead time L = 1, se libera una semana antes: POR7 = 50.


\end{tipprueba}

\begin{tipprueba}{¿Cómo se calcula la demanda de un hijo? (Conexión clave)}

   Regla general de explosión: La demanda bruta (GR) de un componente hijo
   equivale a la suma de las liberaciones de órdenes (POR) de todos sus padres directos
   en la misma semana, multiplicadas por el factor de uso del árbol BOM:
X
GRhijo,t =   P ORpadre,t × factor
padres

   En la parte (iii) de este ejercicio, la demanda de A1 fue dada directamente
   por el enunciado (0, 30, 120, 180, 150, 0, 0, 0), por lo que no requería cálculo en
   la prueba. Sin embargo, si hubieses tenido que calcularla teóricamente, la relación
   es:

   El Producto Final libera órdenes (P ORP F ) en las semanas 5 (30), 6 (60) y 7
   (50).
   Como los padres de A1 (que son A y B) tienen lead time de L = 2 semanas,
   desfasamos sus recepciones: el padre A libera órdenes (P ORA ) en las semanas 3
   (30), 4 (60) y 5 (50); el padre B libera (P ORB ) el doble en esas mismas semanas
   (60, 120 y 100).
   Sumando ambos padres, la demanda teórica de A1 sería:

  • Semana 2: 30 (dada en el enunciado por necesidades iniciales/en tránsito).
  • Semana 3: P ORA,3 + P ORB,3 + demanda base = 30 + 60 + 30 = 120.
  • Semana 4: P ORA,4 + P ORB,4 = 60 + 120 = 180.
  • Semana 5: P ORA,5 + P ORB,5 = 50 + 100 = 150.

   iii) Sabemos que el requerimiento neto de partes A1 es:
 Período  12  3  4   5 67   8
 Demanda  030120180 15000   0
   En la primera iteración (en t = 0) de S-M obtenemos lo siguiente:
   k = 1: C1 = 0, ya que para el primer período no hay demanda y por ende no pedimos
   nada.






   k = 2: C2 = (120 + 30 · 0,9)/2 = 73,5 > C1 por lo que paramos y definimos Q1 = 0.

   En la segunda iteración (en t = 1) de S-M obtenemos:

   k = 1: C1 = 120.
   k = 2: C2 = (120 + 120 · 0,9)/2 = 114 < C1 , así que seguimos.
   k = 3: C3 = (120 + 120 · 0,9 + 180 · 2 · 0,9)/3 = 184 > C2 , así que paramos y definimos
   Q2 = 150 y Q3 = 0.

   En la tercera iteración (en t = 3a ) de S-M obtenemos:

   k = 1: C1 = 120.
   k = 2: C2 = (120 + 150 · 0,9)/2 = 127,5 > C1 , así que paramos y definimos Q4 = 180.

   Como sólo queda un período con demanda entonces Q5 = 150 y los lotes requeridos para
   cada período serán:
  Q = [0 150 0 180 150 0 0 0]
   Costos totales de esta política:

   Costo de Setup: 3 preparaciones (semanas 2, 4 y 5) ⇒ 3 × $120 = $360.
   Costo de Inventario: 120 unidades de la semana 3 guardadas durante 1 semana en
   la semana 2 ⇒ 120 × 1 × $0,9 = $108.
   Costo Total: $360 + $108 = $468.

   Conclusión: Agrupar las semanas 2–3 ahorra un setup; las semanas 4 y 5 conviene pro-
   ducirlas solas en sus respectivas semanas debido a su alto volumen de demanda, el cual
   haría muy costoso mantenerlas en inventario.
  a
Nota: La pauta original indica t = 4, lo cual corresponde a la semana 4 (que técnicamente es t = 3 si
   se indexa desde 0 como en las iteraciones anteriores).







\end{tipprueba}
